% Options for packages loaded elsewhere
\PassOptionsToPackage{unicode}{hyperref}
\PassOptionsToPackage{hyphens}{url}
\documentclass[
]{article}
\usepackage{xcolor}
\usepackage[margin=1in]{geometry}
\usepackage{amsmath,amssymb}
\setcounter{secnumdepth}{-\maxdimen} % remove section numbering
\usepackage{iftex}
\ifPDFTeX
  \usepackage[T1]{fontenc}
  \usepackage[utf8]{inputenc}
  \usepackage{textcomp} % provide euro and other symbols
\else % if luatex or xetex
  \usepackage{unicode-math} % this also loads fontspec
  \defaultfontfeatures{Scale=MatchLowercase}
  \defaultfontfeatures[\rmfamily]{Ligatures=TeX,Scale=1}
\fi
\usepackage{lmodern}
\ifPDFTeX\else
  % xetex/luatex font selection
\fi
% Use upquote if available, for straight quotes in verbatim environments
\IfFileExists{upquote.sty}{\usepackage{upquote}}{}
\IfFileExists{microtype.sty}{% use microtype if available
  \usepackage[]{microtype}
  \UseMicrotypeSet[protrusion]{basicmath} % disable protrusion for tt fonts
}{}
\makeatletter
\@ifundefined{KOMAClassName}{% if non-KOMA class
  \IfFileExists{parskip.sty}{%
    \usepackage{parskip}
  }{% else
    \setlength{\parindent}{0pt}
    \setlength{\parskip}{6pt plus 2pt minus 1pt}}
}{% if KOMA class
  \KOMAoptions{parskip=half}}
\makeatother
\usepackage{graphicx}
\makeatletter
\newsavebox\pandoc@box
\newcommand*\pandocbounded[1]{% scales image to fit in text height/width
  \sbox\pandoc@box{#1}%
  \Gscale@div\@tempa{\textheight}{\dimexpr\ht\pandoc@box+\dp\pandoc@box\relax}%
  \Gscale@div\@tempb{\linewidth}{\wd\pandoc@box}%
  \ifdim\@tempb\p@<\@tempa\p@\let\@tempa\@tempb\fi% select the smaller of both
  \ifdim\@tempa\p@<\p@\scalebox{\@tempa}{\usebox\pandoc@box}%
  \else\usebox{\pandoc@box}%
  \fi%
}
% Set default figure placement to htbp
\def\fps@figure{htbp}
\makeatother
\setlength{\emergencystretch}{3em} % prevent overfull lines
\providecommand{\tightlist}{%
  \setlength{\itemsep}{0pt}\setlength{\parskip}{0pt}}
\usepackage{graphicx}
\usepackage{float}
\usepackage{bookmark}
\IfFileExists{xurl.sty}{\usepackage{xurl}}{} % add URL line breaks if available
\urlstyle{same}
\hypersetup{
  hidelinks,
  pdfcreator={LaTeX via pandoc}}

\author{}
\date{\vspace{-2.5em}}

\begin{document}

\section{Question}\label{question}

Un experimento consiste en medir el alcance horizontal de un proyectil
en funcion del angulo con el que se lanza (respecto a la horizontal). En
la grafica se registran los resultados de 110 lanzamientos realizados
con la misma velocidad inicial.

\includegraphics[width=12cm,height=\textheight,keepaspectratio]{dispersion_proyectil.png}

El comportamiento del alcance respecto al angulo es

\subsection{Answerlist}\label{answerlist}

\begin{itemize}
\tightlist
\item
  lineal y mas disperso cuanto mayor sea el angulo.
\item
  lineal y mas disperso cuanto mayor sea el alcance.
\item
  no lineal y mas disperso cuanto mayor sea el angulo.
\item
  no lineal y mas disperso cuanto mayor sea el alcance.
\end{itemize}

\section{Solution}\label{solution}

Para analizar el comportamiento del alcance respecto al angulo, debemos
observar dos aspectos clave en la grafica de dispersion:

\textbf{1. Tipo de relacion (lineal vs no lineal):}

La relacion entre el angulo y el alcance sigue la ecuacion del
movimiento de proyectiles:

\[R = \frac{v_0^2 \sin(2\theta)}{g}\]

Esta es una funcion senoidal, por lo tanto \textbf{no lineal}. En la
grafica se observa claramente una forma parabolica con un maximo
aproximadamente en \(\theta \approx 0.78\) radianes (\(\approx 45°\)).

\textbf{2. Patron de dispersion:}

Observando la grafica:

\begin{itemize}
\tightlist
\item
  En los extremos (angulos cercanos a 0 y 1.5 radianes): los puntos
  estan mas concentrados
\item
  En la zona central (alcances mayores, alrededor de 10-13 m): los
  puntos muestran mayor variabilidad
\end{itemize}

Esto indica que \textbf{la dispersion es proporcional al alcance}, no al
angulo. Fisicamente, esto tiene sentido porque los errores
experimentales se amplifican en trayectorias mas largas.

\textbf{Conclusion:}

El comportamiento es \textbf{no lineal} (forma parabolica) y \textbf{mas
disperso cuanto mayor sea el alcance} (mayor variabilidad en valores
altos de alcance).

\subsection{Answerlist}\label{answerlist-1}

\begin{itemize}
\tightlist
\item
  \textbf{Falso}. La relacion no es lineal (es parabolica/senoidal) y la
  dispersion no depende del angulo sino del alcance.
\item
  \textbf{Falso}. Aunque la dispersion si aumenta con el alcance, la
  relacion NO es lineal.
\item
  \textbf{Falso}. Si bien la relacion es no lineal, la dispersion NO
  aumenta con el angulo sino con el alcance.
\item
  \textbf{Verdadero}. La relacion es no lineal (forma parabolica) y la
  dispersion aumenta proporcionalmente al alcance.
\end{itemize}

\section{Meta-information}\label{meta-information}

exname:
dispersion\_alcance\_proyectil\_aleatorio\_interpretacion\_representacion\_n2\_v1
extype: schoice exsolution: 0001 exshuffle: TRUE exsection:
Estadistica/Graficas de Dispersion

\end{document}
